\documentclass[10pt,a4paper]{amsart}
\usepackage[margin=19mm]{geometry}
\usepackage{amsmath,amssymb,mathtools}
\usepackage[scr=rsfs,cal=euler]{mathalfa}
\usepackage{xfrac}
\usepackage[unicode,hidelinks,bookmarks=false]{hyperref}
\newcommand{\ie}{i.e.\ }
\newcommand{\cf}{cf.\ }
\newcommand{\resp}{respectively\ }
\newcommand{\Subsection}{\subsection}
\providecommand{\nobreakdash}{\nobreak\hbox{-}}
\setlength{\emergencystretch}{2em}
\multlinegap0pt
\usepackage{mathtools}
\usepackage{amssymb}
\usepackage[shortlabels]{enumitem}
\usepackage{tikz}
\usepackage{bm}
\usepackage{dsfont}
\newtheorem{lemma}{Lemma}
\newtheorem{proposition}{Proposition}
\DeclareMathOperator{\End}{End}
\newcommand{\FF}{\mathbb{F}}
\DeclareMathOperator{\Gal}{Gal}
\DeclareMathOperator{\Ind}{Ind}
\DeclareMathOperator{\Prym}{Prym}
\newcommand{\QQ}{\mathbb{Q}}
\newcommand{\ZZ}{\mathbb{Z}}
\DeclareMathOperator{\one}{\mathds{1}}
\DeclareMathOperator{\wild}{wild}
\title{Wild conductor exponents of curves}
\author{Harry Spencer}
\date{Source: arXiv:2409.12688v3 (CC BY 4.0)}
\begin{document}
\maketitle

\noindent\emph{Source and licence.} Wild conductor exponents of curves. Version arXiv:2409.12688v3 (CC BY 4.0), distributed by arXiv under the Creative Commons Attribution 4.0 International licence, http://creativecommons.org/licenses/by/4.0/, as recorded in the arXiv licence metadata for this exact version and shown on the abstract page https://arxiv.org/abs/2409.12688v3. Full licence notice: \texttt{licenses/arxiv-2409.12688-CC-BY-4.0.txt}.\par\smallskip

\noindent\textit{Editorial scope.} Counted unit: the article's proposition prop:Cq\_main (source0.tex lines 541--546). The article's complete original proof of it is delivered with it (source0.tex lines 548--578, one \texttt{proof} environment of 31 lines). No mathematics is written, repaired or re-derived: the statement and the proof are the article's own lines, in the article's own notation, and no retained line is substituted or re-flowed.  Retained uncounted as the source-local context the proof needs: lines 388--392; lines 433--439; lines 470--473; lines 498--501.  Nothing was omitted from any retained span, so the package carries no omission record.  No retained line contains a citation, so the wrapper carries no bibliography and no bibliography entry.  No retained line contains a figure environment, an image-inclusion command or drawing code, so no image is delivered and the package declares no asset dependency.  Editorial formatting only, in addition: an AMS \texttt{amsart} preamble that restates the article's own declarations and macros used by the retained lines, namely the environments \texttt{lemma}, \texttt{proposition} on separate counters, together with the standard packages those lines need.  The retained environments print in this document as Lemma 1, Proposition 1 in order of appearance, whereas the article prints the same items with the numbering of its own full document.  The complete native source of the article as distributed in the arXiv e-print for this exact version is bundled unchanged as \texttt{sources/165505/source0.tex}.
\begin{lemma}\label{lem:0wild}
    Let $K$ be a finite extension of $\QQ_p$ and fix a prime $\ell\ne p$. Suppose there are abelian varieties $A_1,\hdots,A_n$ and $B_1,\hdots,B_m$ and some wild inertia representation $W$ of dimension $2d$ over $\FF_\ell$ such that
    \[W\oplus\bigoplus_i A_i[\ell]\cong \bigoplus_j B_j[\ell]\]
    as wild inertia representations. If $p>2d+1$, then $n_{\wild}(W)=0$.
\end{lemma}

\begin{proposition}\label{prop:push-pull}
    Let $p$ be a prime and $\pi: X\to B$ a $C_p$-cover of curves over a field $K$ of characteristic different from $p$. Either
    \begin{enumerate}
        \item $\pi$ is unramified, $\ker(\pi_*)/\Prym(\pi)\cong \langle T\rangle$ for some $T\in J_X[p]$ and $\ker\pi^*=\langle P\rangle$ for some $P\in J_B[p]$, or
        \item $\pi$ is ramified, $\ker\pi_*/\Prym(\pi)=1$ and $\pi^*$ is injective.
    \end{enumerate}
\end{proposition}

\begin{lemma}\label{lem:conj_cond}
    Let $X$, $G$ and $\rho$ be as above. For $\sigma \in \Gal(\QQ(\rho)/\QQ)$, we have
    \[ n_{\wild}(X^\rho) = n_{\wild}(X^{\sigma(\rho)}). \]
\end{lemma}

\begin{lemma}\label{lem:curvinduct}
    For $H\le G$, we have
    \[n_{X/H,\wild}=n_{\wild}(X^{\Ind_H^G \one}).\]
\end{lemma}

\begin{proposition}\label{prop:Cq_main}
    Consider a $C_q$ cover of curves $\pi: X\to B$ over a finite extension of $\QQ_p$. Write $R_q$ for the number of ramification points. If $p>\max\{q,(R_q-2)(q-1)+1\}$, then
    \[n_{X,\wild}=q\cdot n_{B,\wild}.\]
    Equivalently, for $\rho$ a non-trivial irreducible representation of $C_q$, we have
    \[ n_{\wild}(X^\rho)=n_{\wild}(X^{\one}). \]
\end{proposition}

\begin{proof}
    By Riemann--Hurwitz, we have $g_X=q\cdot g_B + (R_q-2)(q-1)/2$.

    Now, as an endomorphism on $\Prym(\pi)$, the map $(1-\tau)^{q-1}$ differs from $[q]$ by a unit. This can be seen because $\ZZ[\zeta_q]$ embeds into $\End(\Prym(\pi))$ by sending $\zeta_q$ to $\tau$. Therefore, $\Prym(\pi)[1-\tau]$, the kernel of $1-\tau$ on the Prym, has $\FF_q$-dimension $2(g_X-g_B)/(q-1)=2g_B+R_q-2$ and so $\Prym(\pi)[(1-\tau)^{j}]$ has dimension $j \cdot (2g_B+R_q-2)$.

    Note that, for $R_q\ge2$, we have $\pi^*J_B[q]\le\Prym(\pi)[1-\tau]$ and so 
    \[\Prym(\pi)[1-\tau]\cong V\oplus \ker((1-\tau)\circ\pi^*)[q]\cong V\oplus J_B[q]\] 
    for some $(R_q-2)$-dimensional wild inertia representation $V$ by Maschke's Theorem. In this case, for each $2\le i\le q-1$, we have
    \[ \Prym(\pi)[(1-\tau)^i]/\Prym(\pi)[(1-\tau)^{i-1}] \cong V\oplus J_B[q].\]
    To see this, note that
    \[\Prym(\pi)[(1-\tau)^i]\xrightarrow{(1-\tau)^{i-1}} (1-\tau)^{i-1}(\Prym(\pi))\cap \Prym(\pi)[1-\tau],\]
    is surjective with kernel $\Prym(\pi)[(1-\tau)^{i-1}]$, and so the right-hand side is simply $\Prym(\pi)[1-\tau]$ by counting.

    We decompose $J_X[q]$ as a wild inertia representation:
    \begin{equation}
        J_X[q] \cong J_X[q]/\ker\pi_*[q] \oplus \bigoplus_{i=1}^{q-1} \Prym(\pi)[(1-\tau)^i]/\Prym(\pi)[(1-\tau)^{i-1}].
    \end{equation}
    Using that $J_X[q]/\ker\pi_*[q]\cong\pi_*(J_X[q])$ and that $\pi_*(J_X[q])=J_B[q]$, we have
    \begin{equation}\label{eqn:cond_sum}\tag{$\dagger$}
        J_X[q]\cong J_B[q]^{\oplus q}\oplus V^{\oplus q-1}
    \end{equation}
    as wild inertia representations. Note that $\pi_*(J_X[q])=J_B[q]$ holds because $\ker\pi_*=\Prym(\pi)$ by Proposition \ref{prop:push-pull}, which says that $\ker\pi_*[q]$ has dimension $2g_X-2g_B$, so $\pi_*(J_X[q])$ has dimension $2g_X - (2g_X-2g_B) = 2g_B$. We reach the first conclusion by applying Lemma \ref{lem:0wild} to $V^{\oplus q-1}$, using that $p>\dim V^{\oplus q-1} +1 = (R_q-2)(q-1)+1$. 

    In the case $R_q=0$, there is an analogous identity to \eqref{eqn:cond_sum} with the `junk' terms instead on the left-hand side. One can see this by using the facts that $\pi^*(J_B[q])$ is all of $\ker\pi_*[q]$ and that $\pi_*: J_X[q]\to J_B[q]$ has cokernel of dimension one, then using that $p>q$.

    A monodromy argument shows that we cannot have $R_q=1$. An alternative way to see this is that, assuming $R_q=1$, we would have $\pi^*J_B[q]\le \Prym(\pi)[1-\tau]$, but the left-hand side has dimension $2g_B$ whilst the right has dimension $2g_B-1$. 

    For the second part of the proposition, apply Lemma \ref{lem:curvinduct} to write 
    \[n_{X,\wild}=n_{\wild}(X^{\one})+\sum_{\sigma\in\Gal(\QQ(\zeta_q)/\QQ)}n_{\wild}(X^{\sigma(\rho)}) \text{ and } n_{B,\wild}= n_{\wild}(X^{\one})\]
    and conclude by Lemma \ref{lem:conj_cond}, which says that these conjugate pieces have the same wild conductor exponents.
\end{proof}

\end{document}
